\subsubsection{Why Entanglement Matters}\label{subsubsec:why-entanglement-matters}

Let's answer to the question ``\emph{why entanglement matters}'' step by step. First of all, the \textbf{complexity of quantum states grows exponentially} with the number of qubits, as we have introduced at \autopageref{subsubsec:exponential-explosion}. As consequence, to describe the system, we need $2^n$ complex numbers, where $n$ is the number of qubits. This means that for a system of 50 qubits, we would need $2^{50}$ complex numbers to describe its state, which is an astronomically large number (approximately $1.13 \times 10^{15}$). It could be a \textbf{problem for classical computers}, but \hl{quantum computers can handle this complexity naturally.}

\highspace
The reason why quantum computers can handle this complexity naturally is not that they explicitly store all these amplitudes as a classical computer would. Rather, the quantum system itself physically evolves in this exponentially large state space.

\highspace
Entanglement matters because \textbf{most quantum states cannot be decomposed into independent single-qubit states}. In fact, the set of separable states, i.e., states that can be written as a tensor product of individual qubit states, is \textbf{very small compared to the set of all possible states}. For two qubits, separable states must satisfy:
\begin{equation*}
    c_0 c_3 = c_1 c_2
\end{equation*}
Which is a strong structural constraint. Therefore, \hl{most two-qubit states do not satisfy this condition and are entangled.}

\highspace
\textcolor{Green3}{\faIcon{question-circle} \textbf{But why does this matter?}} Because \textbf{entanglement prevents the system from being compressed into simpler independent components}. The \hl{state must be described globally.} This global structure is difficult to simulate classically, but \hl{quantum mechanics allows us to manipulate it directly through quantum operations.} This is one of the reasons why quantum computers can solve certain problems more efficiently than classical computers.

\highspace
In summary, entanglement is a fundamental feature of quantum mechanics that prevents a multi-qubit system from being decomposed into independent subsystems. As the number of qubits increases, the state space grows exponentially, and entanglement forces a global description of the system. This global structure can be exploited by quantum algorithms to achieve computational advantages over classical approaches.